\subsection{Actual radial arms and separate end caps}
\label{sector:actual-subsection}

The network estimate just proved concerns the law with two full terminal
wires.  We next obtain a bound for actual paths in the blank law.  These
are different assertions: no monotonicity between arbitrary boundary
conditions will be used to pass from one to the other.

Write $H_k(r,R)$ for the straight-wall sector graph of
Section~\ref{sector:radial-setup}, including its two polygonal radial
cuts.  For every fixed
finite integer $k\ge1$, its $k$ successive quadrants are glued
noncyclically on the angular cover, with the two banks kept separate.
In particular, the full-angle case $k=4$ remains a disk with a slit.
Denote its inner and outer vertex arcs by
$A_r=S_r$ and $A_R=S_R$, respectively, and put
\[
 b_k(r,R)=\PP_{H_k(r,R)}
 \bigl(\text{an actual open path joins $A_r$ to $A_R$}\bigr).
\]
Here and below, actual paths use physical edges and cannot use a wire
identification as a step.  A cap partition of an end arc means a
partition drawable by noncrossing auxiliary links in a narrow exterior
cap at that arc.  Inner and outer caps are disjoint.

\begin{proposition}[Blank sector arms]
\label{prop:sector-actual}
For each fixed finite integer $k\ge1$ and each orientation,
\[
 \lim_{M\longrightarrow\infty}\ 
 \sup_{R/r\ge M} b_k(r,R)=0.
\]
If $X$ and $Y$ are separate cap partitions of $A_r$ and $A_R$, with no
other ties, then
\begin{equation}
\label{sector:actual-cap-comparison}
 \PP_{H_k(r,R)/(X,Y)}(A_r\leftrightarrow A_R
       \text{ by an actual path})\le b_k(r,R).
\end{equation}
On the same original edge set, their partition functions satisfy
\begin{equation}
\label{sector:actual-cap-ratio}
 1-b_k(r,R)\le
 \frac{Z_{X,Y}Z}{Z_XZ_Y}\le1.
\end{equation}
\end{proposition}

The comparisons in
\eqref{sector:actual-cap-comparison}--\eqref{sector:actual-cap-ratio}
are immediate applications of Lemmas~\ref{lem:no-hit}
and~\ref{lem:separate-caps}: the sector is a disk and its two radial
cuts are disjoint accessible boundary arcs.  In particular, taking both
caps to be full wires and using the two-wire witness equivalence gives
\[
 Q_k(r,R)
 =\PP_{H_k(r,R)/(A_r,A_R)}
       (A_r\leftrightarrow A_R\text{ by an actual path})
 \le b_k(r,R).
\]
This relation supplies no upper bound on $b_k$ from the smallness of
$Q_k$.

The proof first extracts from any persistent blank arm a sequence of
intermediate bands whose blank crossing probabilities tend to one.
A one-sided collar
comparison then permits the two end wires to be inserted while keeping
a central actual crossing likely.  The exploration and end-extension
lemmas below provide the remaining step: a measured crossing, with one
side uninspected, can be extended across either omitted end gap at a
fixed positive cost.
Applying the two extensions freshly under the full two-wire law would
contradict Proposition~\ref{prop:sector-network}.

\begin{lemma}[A one-sided radial exploration]
\label{sector:frontier}
Consider a sector band with a designated near cut $A$ and far cut $C$;
the designation may run either outward or inward.  Its remaining two
sides are its flanks.  There are two deterministic edge explorations,
one from each flank, with the following properties.
\begin{enumerate}[label=\textup{(\roman*)}]
\item Each exploration detects whether an actual open $A$--$C$ path
exists.  On success it supplies a simple such path $P$, all of whose
edges have been inspected and found open.  Apart from its endpoints,
$P$ avoids the two cuts.
\item All edge interiors strictly on the other side of $P$ from the
starting flank are uninspected.  On deleting inspected closed edges,
the path has an exposed vertex accessible from the starting flank through an
exterior face.
\item On the path-existence event, at least one of the two explorations
has an endpoint on $A$ for which the portion of $A$ between that
endpoint and the uninspected flank has at least half the length of $A$.
\end{enumerate}
Let $E$ be the path-existence event.  For each flank search separately,
let $F_j$, $j=1,2$, be the event that it succeeds and that its own
endpoint leaves at least half of $A$ toward its uninspected flank.
Each $F_j$ is a union of recorded leaves of that search, and
$F_1\cup F_2=E$.  Conditioning on an individual recorded leaf imposes
no conditions on the uninspected edge states beyond the usual induced
FK identifications.
\end{lemma}

\begin{proof}
Keep the order of vertices on both cuts, and treat the cuts as
impassable walls for the complementary search.  Their tangential
physical edges are not queried.  They are unnecessary to detect an
actual radial crossing: trim any such crossing after its last visit
to $A$ and before its first subsequent visit to $C$.  The disk
path/cut alternative for these walls says that either an actual
$A$--$C$ path exists, or a complementary path joins the flanks,
but not both.  One way to verify this deterministic alternative is
to contract each cut to its own terminal in the drawing and then
trim the resulting terminal path as just described.  No probability
law is changed in using this topological test.

Split the complementary exterior into the two flank faces.  Starting
at one of them, explore its entire connected component using
complementary edges, namely crossings of closed physical edges.  The
starting vertex is in an exterior flank collar: entering an adjacent
interior face requires testing the intervening physical boundary edge
and finding it closed.  In particular, adjacent interior faces are
not seeded for free.  No passage out through $A$ or $C$ is allowed.
One implementation keeps a
fixed ordered queue of reached complementary faces and tests every
previously untested physical edge incident to the next face.  A closed
edge adds its opposite face to the queue.  An open edge does not.
Stop on reaching the other flank, or on exhausting the queue.  Thus the
history is generated by successive tests of individual edge states.

If the other flank is reached, the complementary path blocks every
actual $A$--$C$ path.  Otherwise set every uninspected physical edge
temporarily closed.  The reached complementary component does not
enlarge: every edge leaving it was inspected and found open.  There
is still no complementary flank-to-flank crossing, so the disk
path/cut alternative supplies an open radial crossing in this
temporary configuration.  All its physical edges were consequently
inspected and found open.  Choose a simple one by fixed deterministic
tie breaks, and retain the portion after its last near-cut visit and
before its first subsequent far-cut visit.  This defines $P$ and
proves~\textup{(i)}.  Temporary closure is only a device for selecting
the witness, not additional conditioning of the actual configuration.

Every reached face lies on the starting side of any open radial
blocker.  In particular no tested edge interior lies strictly on the
opposite side of the chosen frontier: an edge is tested only from an
incident reached face.  The closed edges crossed during the search
give a chain of face expansions from the starting flank.  Deleting
them exposes the appropriate side of $P$ to that exterior face.
These observations prove~\textup{(ii)}.  If $P$ touches a flank, use
separate exterior boundary slots, or equivalently slightly prolong
that flank outward, to distinguish its two incident sides.  Only the
component accessible from the opposite flank is relevant; any further
pieces separated by a boundary touch attach at $P$ after contraction.

Orient the near-cut arc from the first flank to the second.  Among
simple actual radial paths, with internal visits to the cuts removed,
consider the possible near endpoints.  The search from the first
flank has a leftmost such endpoint.  To check this directly, let $C_0$
be any such blocker.  The exploration cannot reach a face strictly
beyond $C_0$.  An inward edge with near endpoint strictly beyond the
endpoint of $C_0$ has its entire edge interior on the blocked side;
neither incident complementary face is reached, so it cannot have
been inspected.  Hence the near endpoint of the selected
inspected-open path cannot be farther along the arc than the
endpoint of $C_0$.  The reversed search similarly gives
a rightmost endpoint.  If their arc coordinates are $x_-$ and $x_+$
in an arc of length $\ell$, then $x_-\le x_+$ and
\[
 \max\{\ell-x_-,x_+\}\ge\ell/2.
\]
This gives~\textup{(iii)} and $F_1\cup F_2=E$.  For each search, its
frontier and its own clearance test are functions of that search's
recorded leaf.  An adaptive edge exploration with deterministic
next-edge choices leaves the ordinary FK domain Markov law on untested
edges, proving the final assertion.  No joint history of the two
searches is used.
\end{proof}

\begin{lemma}[Extension across one radial end gap]
\label{sector:end-extension}
Fix an integer $L\ge2$.  Let $R/r$ be sufficiently large and $r$
sufficiently large, and set
\[
 a=2Lr,\qquad s=\lfloor R/(2L)\rfloor.
\]
In $H_k(r,R)$ give $A_r$ and $A_R$ their full separate wires.  There
is a constant $c>0$, depending only on $q,L,k$, such that
\begin{align}
\label{sector:inner-extension}
 \PP(\text{actual $r$--$s$ crossing})
 &\ge \frac c2\,
      \PP(\text{actual $a$--$s$ crossing}),\\
\label{sector:outer-extension}
 \PP(\text{actual $r$--$R$ crossing})
 &\ge \frac c2\,
      \PP(\text{actual $r$--$s$ crossing}).
\end{align}
Every crossing on these lines is required to stay between its
indicated radial cuts.  The searches establishing the two estimates
are applied separately under this same original law.
\end{lemma}

\begin{proof}
For the first estimate the explored band is $[a,s]$, its near cut is
$a$, and the targeted end is $r$.  For the second it is $[r,s]$,
its near cut is $s$, and the targeted end is $R$.  In either case
write $E$ for the actual crossing event in the explored band.
In the first instance the opposite original end wire at $R$ lies
beyond the retained band.  In the second, the opposite original end
wire at $r$ meets the measured path on the retained far cut.  This is
the only difference in the removal of end identifications below.
Apply Lemma~\ref{sector:frontier} from both flanks.  On a favorable
leaf $\omega$, namely a successful search with its own clearance in that
lemma, call the uninspected flank the right flank and its measured
open path $P$.  Let $G_\omega$ be the conditional graph obtained from
the full $H_k(r,R)$ by deleting the inspected closed edges,
contracting the inspected open edges, and imposing the two original
separate end wires.  The remaining edges in this graph have exactly
the FK law conditional on $\omega$.  Write $T$ for the targeted end
terminal in $G_\omega$; all vertices of $P$ already belong to one
terminal there.  Thus the probability of contact in the full network
conditional on $\omega$ is $\PP_{G_\omega}(P\leftrightarrow T)$.

Attach to the ordinary explored band a closed lattice rectangle running
along the right flank across the missing end gap.  Its transverse width is
\begin{equation}
\label{sector:extension-width}
 w=\left\lfloor
       \frac{\min\{\text{target radius},\text{near radius}\}}8
       \right\rfloor.
\end{equation}
The rectangle lies in the quadrant adjacent to that flank.  Its
opening on the near cut lies strictly to the right of the endpoint
of $P$: the clearance is at least half the whole near arc, whereas
$w$ is at most one eighth of its radius.  Its aspect ratio is
bounded in terms of $L$; $w$ is positive at the scales under
consideration.  Let $H_*$ be the ordinary explored band together with
this rectangle, and let $B_0$ be the rectangle's end interval on the
targeted radial cut.

Define the reference graph
\[
 R_\omega=H_* /(P,B_0).
\]
It has all ordinary edges of the explored band and attached rectangle,
with the vertices of $P$ and $B_0$ wired separately and no other
identifications.  In particular it has none of the other measured
deletions or contractions in $G_\omega$.  We claim
\begin{equation}
\label{sector:conditional-extension}
 \PP_{G_\omega}(P\leftrightarrow T)
 \ge \PP_{R_\omega}(P\leftrightarrow B_0)
 =g_{H_*}(P,B_0).
\end{equation}

We compare these two graphs through changes that cannot increase
the tested probability in $G_\omega$.  For these comparisons, represent
its fixed ties by auxiliary links of finite odds and take their
contraction limits afterward.
Shrinking the targeted end wire to $B_0$ does not increase the tested connection
probability by adjacency; its removed vertices remain ordinary
vertices of the graph.  The deleted closed edges expose $P$ to the
starting-flank exterior.  Draw the terminal test edge through that
accessible region and the exterior to an exterior representative of
$B_0$.  Keep auxiliary cap chains narrow.

Every edge outside $H_*$ can now be removed by cofacial face expansion.
On the target side, the part of the end gap omitted by the rectangle
opens through the remaining target boundary or through a flank; the
restricted $B_0$ cap seals neither access.  On the far side, the
omitted part opens through the far boundary and the flanks.  Thus
each unwanted edge has a midpoint accessible from a test-edge face
across only unwanted edges.  If the opposite full end wire is beyond
the retained band, remove its auxiliary cap links first from that
same exterior.  For each finite comparison, cofacial Rayleigh says
that the ensuing deletion does not increase the network connection
probability.  Fixed open ties can equivalently be assigned odds
tending to infinity and then contracted, so this argument applies
also in their presence.

Delete the vertices of $P$ from $H_*$ and take the component accessible
from the rectangle.  Define $U_\omega$ by adjoining to this component
its incident edges into $P$ and the terminal vertex sets $P,B_0$.
Every edge interior of $U_\omega$ outside $P$ is uninspected and has its
ordinary odds.  Now contract $P$; its own edges become loops and may
be omitted.
The searched side, and any pieces cut off from the rectangle by
contacts of $P$ with a flank, attach to $U_\omega$ only at this
contracted vertex.  In the second application the retained opposite
end wire also meets $P$; remove its additional
identifications first.  This again does not increase the tested probability by
adjacency at the $P$ terminal.

Now compare the reduced conditional graph with $R_\omega$.  They
have exactly the same ordinary graph on $U_\omega$, with $P$ and
$B_0$ as separate terminals.  In either graph every remaining piece
outside $U_\omega$ attaches to it only at $P$.  An FK graph attached
through one vertex contributes a partition-function factor independent
of the other side's states, so those pieces cancel from the terminal
connection marginal.  The marginal is therefore the one in
$R_\omega$.  Its tested probability is therefore at most that in
$G_\omega$, proving~\eqref{sector:conditional-extension}.

It remains to bound its right-hand side uniformly in the history.
Choose an interval $J$ of fixed positive relative length on the
left band flank, with radius comparable to the near radius and away
from both radial cuts.  In the first application retain temporarily
only the rectangle and the subband $[a,4a]$, choosing $J$ there.  In
the second retain the rectangle and
$[\lfloor s/4\rfloor,s]$, again choosing $J$ in its interior radial
range.  The assumed large ratio ensures that these subbands fit.

There is a boundary route from $B_0$ along the right flank, around
the truncation rim opposite the rectangle opening, and back along
the left flank to $J$.  In the first case it uses the outer rim; in
the second it uses the inner rim.  Along the rectangle, its inward
clearance is a fixed fraction of $w$, and along the sector it is a
fixed fraction of the near radius.  Their ratio is bounded in terms
of $L$.  Straight portions and corner neighborhoods therefore admit
the hypotheses of Lemma~\ref{sector:blank-boundary-chain}, with
constants depending only on $L,k$.  At the inner-rim corners use
the clipped-square construction in that lemma.  At split banks use
their own sides of the abstract disk, never a neighborhood crossing
the slit.  It follows that the pure $B_0$--$J$ connection probability
in this truncated shape is at least $c(q,L,k)>0$.

Extend to $H_*$ cofacially.  The further band pieces open to the
exterior beyond the truncated rim; the two terminal caps are on
$B_0,J$ and do not obstruct this access.  Consequently
\[
 g_{H_*}(B_0,J)\ge c(q,L,k).
\]
Every graph route from the attached rectangle to $J$ meets $P$,
because $P$ is a radial separator and the rectangle opening is on
its right.  Add $P$ to the $J$ terminal by adjacency.  Once $P$ is
contracted, all the remaining pieces on the $J$ side attach only at
that terminal and factor out.  Wiring their $J$ vertices to it
does not change the other side's marginal law.  Therefore
\[
 g_{H_*}(P,B_0)
   =g_{H_*}(P\cup J,B_0)
   \ge g_{H_*}(J,B_0)\ge c(q,L,k).
\]
Together with~\eqref{sector:conditional-extension}, this is the
required lower bound given each favorable exploration history.

A successful network connection here supplies the indicated actual
extension.  First expand every contracted measured-open component
along its already measured physical edges.  The only remaining
abstract jumps are in the two original end wires.  Starting from the
target terminal, trim its path after its last use of the target wire.
If it reaches the opposite end wire or leaves the explored band past
the far cut before reaching $P$, its physical portion has already
reached that cut.  Otherwise
it reaches the measured physical path $P$, which continues to the
far cut.  Erasing loops gives an actual crossing, stopped at its
first visit to the far cut.  A target-wire jump only chooses the
starting vertex of this physical path.

For the last probability step, let $F_1,F_2$ be the favorable history
events for the two flank explorations.  They are both contained in
$E$, and Lemma~\ref{sector:frontier} gives $F_1\cup F_2=E$.
Apply the conditional estimate separately to each exploration, not
to a joint unexplored law.  If $E'$ is the extended actual-crossing
event, then
\[
 \PP(E')\ge c\PP(F_i)\quad(i=1,2),\qquad
 \PP(E')\ge c\max_i\PP(F_i)\ge\frac c2\PP(E).
\]
This proves both asserted inequalities.  In particular the second
one is a fresh application under the original law, not an application
under a law conditioned on the outcome of the first search.
\end{proof}

\begin{proof}[Proof of Proposition~\ref{prop:sector-actual}]
Only the qualitative blank decay remains to be shown.  Fix $k$ and
an orientation.  If it fails, there is $\epsilon>0$ and sectors with
arbitrarily large $R/r$ for which $b_k(r,R)\ge\epsilon$.

Choose such a sector with $R/r$ so large that it contains $n$
mutually disjoint radial bands, each of ratio at least $n$, and
each of inner radius at least $n$.  The bands may be separated by
additional radial gaps.  This is possible for every $n$: for
example, choose the original ratio larger than $n^{4n}$ and use
successive geometrically spaced integer radii, increasing the
ratio slightly to absorb roundoffs.  Removing one closed band
leaves an inner piece and an outer piece with no physical route
between them.  Conditional on all edges outside that band, each
piece induces a noncrossing partition only on its own radial cut.
There are no lateral wires and no through ties.  By
\eqref{sector:actual-cap-comparison}, the conditional crossing
probability in that band is bounded by its own blank probability.

A full-sector crossing requires crossings of all the bands.
Conditioning outside one band at a time therefore gives
\[
 \epsilon\le b_k(r,R)
      \le \prod_{i=1}^n b_k(r_i,R_i).
\]
No independence is asserted here.  At least one factor is at least
$\epsilon^{1/n}$.  Selecting such a band for every $n$ produces a
new sequence, renamed $(r,R)$, with
\begin{equation}
\label{sector:blank-near-one}
 r\longrightarrow\infty,\qquad R/r\longrightarrow\infty,
 \qquad b_k(r,R)\longrightarrow1.
\end{equation}

Choose a fixed integer $L\ge2$ sufficiently large that
Proposition~\ref{prop:sector-network} gives
\begin{equation}
\label{sector:collar-pure-small}
 Q_k(x,y)\le q-\varepsilon_0
 \qquad\text{whenever }y/x\ge L
\end{equation}
for some fixed $\varepsilon_0>0$.  Enlarge $L$ once if necessary
for integer cut conventions.  Consider the actual crossing of the
central band
\[
 [a,s]=[2Lr,\lfloor R/(2L)\rfloor].
\]
Its probability tends to one in the blank law by
\eqref{sector:blank-near-one}, since a full crossing contains such
a central crossing.

Insert the full inner end wire, using its collar $[r,Lr]$.
Conditional on the exterior of this collar, the exterior gives
only an opposite-cut cap partition.  Lemma~\ref{lem:one-sided-density}
and~\eqref{sector:collar-pure-small} bound the density and its
reciprocal on that exterior by a constant $C_0<\infty$ independent
of $r,R$.  For example, the logarithmic variation is at most
\[
 \log\left(1+\frac{(t-1)(q-\varepsilon_0)}
                         {1-t(q-\varepsilon_0)}\right),
\]
whose denominator is positive.  Now insert the full outer end wire,
using $[\lfloor R/L\rfloor,R]$.  The same estimate applies:
the already inserted inner wire is part of the exterior and still
induces only a noncrossing cap partition on the collar's opposite
cut.  The central event is exterior to both collars.  Its failure
probability after both insertions is at most $C_0^2$ times its
blank failure probability, and hence tends to zero.

Apply the two estimates of Lemma~\ref{sector:end-extension} under
this law with full separate end wires.  Their constants are fixed,
because $q,L,k$ are fixed.  They give
\[
 \liminf_{R/r\to\infty}
 \PP_{H_k(r,R)/(A_r,A_R)}
       (\text{actual $r$--$R$ crossing})
 \ge \frac{c(q,L,k)^2}{4}>0
\]
along the sequence~\eqref{sector:blank-near-one}.  With only these
two separate full wires, terminal network connection is equivalent
to existence of an actual path between the two end arcs: take a
simple terminal path and remove its initial and final wire portions.
The displayed probability is therefore $Q_k(r,R)$, contradicting
Proposition~\ref{prop:sector-network}.  This proves the blank decay,
and the end-cap comparisons were established at the start.
\end{proof}

\begin{corollary}[Power-law sector bound]
\label{cor:sector:power}
For each fixed finite integer $k\ge1$ there are constants $C<\infty$
and $\eta>0$, depending on $q,k$, such that for all integer
$1\le r<R$,
\begin{equation}
\label{sector:power-arm}
 b_k(r,R)\le C(r/R)^\eta.
\end{equation}
The same upper bound holds for actual crossings under arbitrary
separate inner and outer cap partitions, with no other ties.
Their partition-function ratio in
\eqref{sector:actual-cap-ratio} lies between
$1-C(r/R)^\eta$ and $1$.
\end{corollary}

\begin{proof}
Choose an integer $M\ge2$ so large that the supremum in
Proposition~\ref{prop:sector-actual} at ratio $M$ is at most some
fixed $\theta<1$.  Between $r$ and $R$ place the bands
\[
 [(2M)^j r,\ M(2M)^j r],
\]
for all consecutive integers $j\ge0$ for which they fit.  Their
number $N$ is at least
$\log(R/r)/\log(2M)-1$, unless this lower bound is negative, in
which case the eventual estimate follows by increasing $C$.
The intervening gaps make the edge sets disjoint.  Conditioning
outside each band induces only separate end caps, so the same
iteration as above gives a crossing probability at most $\theta^N$.
This argument also applies with prescribed separate cap partitions
at the two original ends, because neither partition joins the
inner and outer components left by removing a band.  Taking
\[
 \eta=\frac{-\log\theta}{\log(2M)}
\]
and absorbing the one-band discrepancy into $C$ proves the first
two assertions.  The last follows from
\eqref{sector:actual-cap-ratio}.
\end{proof}

\begin{remark}
All constants above are for fixed $q$ and fixed finite $k$.  The
proof uses a number of boundary portions bounded in terms of $k$,
and always works on the disk with its two distinct banks.  It gives
no uniformity as $k$ increases, and none of its assertions compares
arbitrary through partitions or lateral wires.
\end{remark}
